\section{Conclusion and Future Work}\label{sec:conclusion}

The 50\% GER target is achievable for about half of the States under a continuation of recent trends, but not for the States that account for most of India's young population. College density is the most robust correlate of GER and the lever with the largest projected effect. PTR and GPI are important for quality and equity, but they do not carry independent explanatory power for the aggregate in this cross-section. The population-weighted GER is projected at 43.5\% in 2035-36 under the reference scenario, 6.5 points short of the target, with 68.6\% of the 18--23 population in States below 50\%.

The policy message is twofold. First, the target is off track in the populous low-density States, and a national strategy that focuses on the already-high-GER States will not deliver it. Second, capacity expansion must be paired with teacher supply and equity policies, because the quality and parity levers in this model are too weak to close the gap on their own. A credible 2035 strategy needs a differentiated, State-specific mix of expansion, quality and equity, monitored with uncertainty rather than a single point forecast.

Future work should extend the panel beyond 2021-22, add a second year of AISHE data or a district-level check for a few States, and use institution-level data to reduce aggregation bias. It should also test whether the density association survives a credible identification strategy such as a policy-induced expansion programme, and model female GER and field-of-study segregation directly rather than treating gender parity as a regressor for the aggregate.
